% ============================================================
% CAPITOLUL 6 — EVALUAREA PERFORMANȚEI ȘI SCALABILITĂȚII
% ============================================================

\chapter{Evaluarea Performanței și Scalabilității}
\label{ch:evaluare}

\section{Metodologia de benchmarking}
\label{sec:metodologie_benchmark}

\subsection{Configurații hardware utilizate}
\label{subsec:hardware}

\todosec{Tabel cu configurațiile hardware: laptop (Intel i5, 8GB RAM),
workstation (AMD Ryzen 9, 32GB RAM), server cloud (AWS EC2 c5.4xlarge).
Sistemul de operare, versiunile limbajelor, versiunile solver-elor.}

\subsection{Metrici urmărite}
\label{subsec:metrici}

\todosec{Metrici: (1) Timp de execuție wall-clock (mediana și IQR peste 30 rulări),
(2) CPU time (user + system), (3) Peak memory usage (RSS), (4) Acuratețe numerică
(eroarea față de soluția de referință DOPRI8 cu toleranță $10^{-12}$).}

\subsection{Protocoale de testare și reproductibilitate}
\label{subsec:protocoale}

\todosec{Protocol: warm-up runs, izolare procese, seed fix, versioning
al codului de benchmark. Toate configurațiile rulate de 30 de ori, cu
colectarea automatizată a rezultatelor.}

\section{Scalabilitate în funcție de dimensiunea setului de date}
\label{sec:scalabilitate_date}

\todosec{Experimente cu $N \in \{10^3, 10^4, 10^5, 10^6, 10^7\}$.
Grafice log-log ale timpului de execuție față de $N$ pentru fiecare limbaj.
Identificarea complexității empirice.}

\section{Comparație timp de execuție — R vs Julia vs Python}
\label{sec:comparatie_timp}

\todosec{Tabel comparativ și grafice boxplot. Analiza statistică:
test Kruskal-Wallis, comparații post-hoc Dunn pentru detectarea diferențelor
semnificative. Discutarea overhead-ului de start-up al fiecărui limbaj
față de timpul de calcul pur.}

\section{Consum de memorie și resurse}
\label{sec:consum_memorie}

\todosec{Profile de memorie pentru fiecare implementare. Grafice de utilizare
a memoriei în timp. Identificarea cauzelor alocărilor excesive.}

\section{Acuratețea numerică a rezultatelor}
\label{sec:acuratete}

\todosec{Compararea soluțiilor numerice cu soluția de referință.
Eroarea $L^\infty$ și $L^2$ pentru $S(t)$, $I(t)$, $R(t)$.
Verificarea conservării $S + I + R = N$ (eroarea de conservare).}

\section{Performanța bazei de date distribuite}
\label{sec:performanta_bd}

\subsection{Latență la scriere și citire}
\label{subsec:latenta_bd}

\todosec{Măsurarea latenței P50/P95/P99 pentru operații INSERT și SELECT
pe tabelele epidemiologice. Comparație cu PostgreSQL standalone.}

\subsection{Comportament sub load — teste de stress}
\label{subsec:stress_bd}

\todosec{Teste de load cu \texttt{pgbench} și/sau \texttt{cassandra-stress}.
Throughput (tranzacții/secundă) față de numărul de clienți concurenți.
Comportamentul clusterului la oprirea unuia din cele 3 noduri YugabyteDB.}

\section{Performanța API-ului Go}
\label{sec:performanta_api}

\todosec{Benchmarking API cu \texttt{wrk} sau \texttt{hey}: RPS (requests/second),
latență P95/P99, comportament sub concurență (10, 50, 100, 500 clienți simultani).}

\section{Analiză comparativă sintetică}
\label{sec:analiza_sintetica}

\todosec{Tabel de scoruri normalizate per metrică și per limbaj.
Radar chart (spider chart) cu performanța relativă.
Recomandări pentru scenarii de utilizare: volum mic vs. mare, precizie vs. viteză.}
